\chapter{Concepts you need}\label{ch:concepts}

This chapter is for a reader with little CFD practice. It gives each concept that the
workflow uses, why the workflow needs it, and where to learn more. Do the OpenFOAM
tutorials of \cref{tab:tutorials} before your first run of this case.

\section{An OpenFOAM case}

An OpenFOAM case is a directory with three parts:

\begin{tabularx}{\linewidth}{@{}P{0.18\linewidth} L@{}}
\toprule
Directory & Contents \\
\midrule
\file{0/} & The initial and boundary conditions of each field ($p$, $\mathbf{U}$, $T$, \ldots). The templates keep them in \file{0.orig/} and copy them to \file{0/} at the start. \\
\file{constant/} & The mesh (\file{polyMesh/}) and the physical models (gas, turbulence). \\
\file{system/} & The numerical settings: time control (\file{controlDict}), discretization (\file{fvSchemes}), linear solvers (\file{fvSolution}), parallel decomposition. \\
\bottomrule
\end{tabularx}

\medskip
Each file is a dictionary of \texttt{keyword value;} entries. \cmd{foamDictionary}
reads and writes them from the shell. The OpenFOAM user guide explains the structure.

\section{Finite volumes and the mesh}

The finite-volume method divides the domain into cells and solves the conservation of
mass, momentum and energy for each cell: what enters through the faces equals what
leaves plus what accumulates \citep{versteeg2007,moukalled2016}. The value at a face
comes from the values at the cell centres with an interpolation scheme. The accuracy
thus depends on the size and the shape of the cells where the flow changes fast:
shocks, the boundary layer, the wake. A refined mesh reduces the discretization error;
the mesh-convergence study of \cref{sec:gci} measures it.

Mesh quality has three common measures. \emph{Non-orthogonality} is the angle between a
face normal and the line between the two cell centres. \emph{Skewness} is the distance
of that line from the face centre. A \emph{negative volume} is an inverted cell, which
makes the mesh unusable.

\section{Compressible flow regimes}

The Mach number $M = U/a$ sets the physics and the numerics \citep{anderson2003}:

\begin{tabularx}{\linewidth}{@{}l L@{}}
\toprule
Regime & What happens \\
\midrule
$M < 0.3$ & Density changes are small. \\
$0.3 < M < 0.8$ & Compressible, without shocks. \\
$0.8 < M < 1.2$ & Transonic: local shocks, a fast rise of the drag. \\
$M > 1.2$ & Supersonic: a bow shock at the nose, oblique shocks at the fins, an expansion and a wake at the base. Information travels only downstream, inside the Mach cone. \\
\bottomrule
\end{tabularx}

\medskip
This is why the workflow uses two solvers (\cref{sec:two-solvers}) and different outlet
conditions (\cref{sec:bc}).

\section{Turbulence, RANS and \texorpdfstring{$y^+$}{y+}}

The flow around a rocket is turbulent. A RANS model solves for the mean flow and
models the effect of the turbulent fluctuations with an eddy viscosity $\nu_t$
\citep{wilcox2006,pope2000}. The $k$--$\omega$ SST model of this workflow solves two
equations, for the turbulent kinetic energy $k$ and its specific dissipation rate
$\omega$.

Near a wall the velocity changes over a very small distance. $y^+ = y\,u_\tau/\nu$ is
the distance of the first cell centre from the wall in wall units. Below $y^+ \approx 5$
the cell is in the viscous sublayer; above $y^+ \approx 30$, in the log layer; between
them, in the buffer layer. A wall function gives the wall shear stress from the first
cell when the mesh does not resolve the sublayer (\cref{sec:wall-treatment}).

\section{Steady solution by pseudo-time marching}

The workflow wants the steady flow. The solvers march in a pseudo-time from an initial
field until the solution stops changing (\cref{sec:lts}). Each iteration reduces the
\emph{residuals}, the error of the discrete equations, and the force coefficients become
flat. The run stops when the mean of $C_D$ over a window is constant
(\cref{sec:stop-criterion}). The \emph{Courant number} $Co = U\,\Delta t/\Delta x$ limits
the step of each cell: a large step makes the march unstable.

\section{Force coefficients}

The force on the rocket, divided by the dynamic pressure $\tfrac12\rho_\infty U_\infty^2$
and a reference area, gives the coefficients. $C_D$ (drag) and $C_L$ (lift) are along
and normal to the free stream; $C_A$ (axial) and $C_N$ (normal) are along and normal to
the body axis. $C_m$ is the moment coefficient about a reference point. The centre of
pressure is where the normal force acts; a rocket is stable if it is behind the centre
of gravity.

\section{Learning path}

\begin{table}[htbp]
\centering
\caption{OpenFOAM tutorials to do before this case, in this sequence. Paths are relative to \texttt{\$FOAM\_TUTORIALS}.}\label{tab:tutorials}
\begin{tabularx}{\linewidth}{@{}P{0.46\linewidth} L@{}}
\toprule
Tutorial & What it teaches \\
\midrule
\file{incompressible/icoFoam/cavity/cavity} & The case structure, a run, ParaView. \\
\file{incompressible/simpleFoam/pitzDaily} & A RANS model, wall functions, convergence of a steady solution. \\
\file{compressible/rhoPimpleFoam/RAS/aerofoilNACA0012} & The pressure-based compressible solver with RANS on an external flow. \\
\file{compressible/rhoCentralFoam/forwardStep} & The density-based solver and a supersonic flow with shocks. \\
\file{compressible/rhoCentralFoam/wedge15Ma5} & An oblique shock with an exact solution (\cref{sec:supersonic-validation}). \\
\bottomrule
\end{tabularx}
\end{table}

To copy and run a tutorial:
\begin{shell}
cp -r $FOAM_TUTORIALS/incompressible/icoFoam/cavity/cavity ~/cavity
cd ~/cavity && ./Allrun   # or: blockMesh && icoFoam
\end{shell}

For the theory, read the finite-volume method in \citet{versteeg2007} or, with OpenFOAM
examples, \citet{moukalled2016}; compressible flow in \citet{anderson2003}; turbulence
modelling in \citet{wilcox2006}.
